% !PRINTED 661
\weberpage{661}
%<<CONT
worin $\xi$ die in \S.\,172 näher bestimmte Zahlenreihe durchläuft.
Es möge $\mathfrak{p}_\xi^{k_\xi}$ die in $\varphi$ enthaltene Potenz von
$\mathfrak{p}_\xi$ sein, so dass, wenn wir
\[
\varphi = \chi \prod^{\xi} \mathfrak{p}_\xi^{k_\xi}
\]
setzen, $\chi$ theilerfremd zu $p$ und mit keinen anderen Primzahlen verwandt
ist (s. den Schluss von \S.\,173), als $\varphi$ selbst. Daraus ergiebt sich
\weq{5}{\varphi_n = \chi_n \prod^{\xi} \mathfrak{p}_{n\xi}^{k_\xi},}
und hierin lassen wir $n$ abermals die Reihe der Zahlen $\xi$ durchlaufen.
Setzen wir noch
\[
\psi = \prod^{\xi} \chi_\xi,
\]
so ist auch $\psi$ theilerfremd zu $p$ und mit keinen anderen Primzahlen
 verwandt als $\varphi$, und wir erhalten aus (5) mit Benutzung von (4):
\weq{6}{\prod^{\xi} \varphi_\xi = \psi\, p^{\Sigma k_\xi},}
und daraus ergiebt sich, dass das Functional $\psi$ denselben beiden
Bedingungen (1), (2) genügt, wie $\varphi$.

Es zeigt sich ferner, dass, wenn die $\varphi_\xi$ in der Hauptclasse
liegen, auch $\psi$ in der Hauptclasse enthalten ist.

Wenn aber $p-1$ nicht durch $q$ theilbar ist, so können wir
auch das Umgekehrte schliessen.

Denn nach (2) ist $\varphi_\xi$ äquivalent mit $\varphi^\xi$ und folglich ist:
\[
\begin{gathered}
\prod^{\xi} \varphi_\xi \ \text{äquivalent}\ \varphi^{\Sigma\xi}, \\
\text{äquivalent}\ \psi,
\end{gathered}
\]
und wenn $\psi$ mit einer Zahl associirt ist, so gilt dasselbe
von $\varphi^{\Sigma\xi}$. Nach dem Satze \S.\,172 ist aber, wenn $p-1$ nicht
durch $q$ theilbar ist, $\Sigma \xi$ auch nicht durch $q$ theilbar, und folglich
 ist auch $\varphi$ selbst in der Hauptclasse enthalten.

Diese Betrachtung können wir nun, wenn $\psi$ noch mit
weiteren Primzahlen, die nach dem Modul $q$ nicht mit 1 congruent
 sind, verwandt ist, wiederholen, und kommen so zu dem
Satze:

\gesperrt{Der Satz a) ist allgemein bewiesen, wenn er
unter der Voraussetzung bewiesen werden kann,
dass $\varphi$ nur mit solchen Primzahlen verwandt ist,
die nach dem Modul $q$ mit 1 congruent sind.}
